\section{Introdução}\label{into}

% Abstract—In the current DVB generation, satellite terminals
% are expected to be interactive and capable of transmission in the
% return channel with satisfying quality. Considering the bursty
% nature of their traffic and the long propagation delay, the use of a
% Random Access technique is a viable solution for such a Medium
% Access Control (MAC) scenario. In this paper Random Access
% communication design in DVB-RCS2 is considered with particular regard to the recently introduced Contention Resolution
% Diversity Slotted Aloha (CRDSA) technique. The paper presents
% a model for design and tackles some issues on performance
% evaluation of the system by giving intuitive and effective tools.
% Moreover, dynamic control procedures able to avoid congestion at
% the gateway are introduced. Results show the advantages brought
% by CRDSA to DV

Na geração atual de DVB (Digital Video Broadcast), é esperado que terminais do satélite sejam iterativos e capazes de transmitir dados com uma alta qualidade. Considerando o grande necessidade de tráfego e uma longa propagação de delay , uma técnica de acesso randômica é a solução para tal problema, esse acesso randômico é feito em DVB-RCS2 \cite{DBLP:journals/corr/MeloniM15c}

Todos os padrões DVB são desenvolvidos ou endossados pelo DVB Project, que é uma associação da indústria de mais de 200 fornecedores e usuários de equipamentos e software comprometidos com padrões técnicos abertos para serviços de dados e televisão digital. Embora os padrões tradicionais de transmissão de televisão digital permaneçam importantes, os padrões DVB de maior interesse e impacto futuro são aqueles que fornecem comunicações digitais de alta velocidade para IP (Internet Protocol) e acesso à Internet via satélite, como DVB-RCS2 e DVB-S2.

O DVB-RCS, e especialmente a sua segunda geração (DVB-RCS2), são de grande importância neste aspecto, suportando comunicações de satélite interactivas de alta velocidade (bidireccionais) até 150 Mbps ou mais para transferências (Rx) e 75 Mbps ou mais para uploads (Tx) quando fornecido com capacidade de transponder adequada e configurações VSAT. Portanto, velocidades muito mais altas em conexões sem fio para assinantes são possíveis do que com redes móveis terrestres 4G atualmente.
% \IEEEPARstart{C}{I}rcuitos de duas portas ou quadripolos são topologias de circuitos caracterizados pelos seus parâmetros de impedância e permitem simplificar a associação destes por espécies de módulos. As aplicações são infinitas, como filtros, amplificadores de potência, casamento de impedâncias para eletrônica de Radio Frequência e muitas outras funcionalidade oriundas dos circuitos eletrônicos. 
  
% A determinação prática dos parâmetros de impedância segue a referência da fórmula matricial \cite{b1} que relaciona os parâmetros de impedância com as tensões e correntes do quadripolo.

% \begin{equation}\label{eq:1}
% \begin{bmatrix}
% V_1\\
% V_2
% \end{bmatrix}
% =
% \begin{bmatrix}
% Z_{11} & Z_{12}\\
% Z_{21} & Z_{22}\\
% \end{bmatrix}
% .
% \begin{bmatrix}
% I_1\\
% I_2
% \end{bmatrix}
% \end{equation}

% Desta maneira, é possível representar cada impedância separadamente.

% \begin{multicols}{2}
%   \begin{equation}\label{eq:2}
%     Z_{11} = \frac{V_1}{I_1}\Bigr|_{\substack{I_2=0}}
%   \end{equation}\break
%   \begin{equation}\label{eq:3}
%     Z_{12} = \frac{V_1}{I_2}\Bigr|_{\substack{I_1=0}}
%   \end{equation}
% \end{multicols}

% \begin{multicols}{2}
%   \begin{equation}\label{eq:4}
%     Z_{21} = \frac{V_2}{I_1}\Bigr|_{\substack{I_2=0}}
%   \end{equation}\break
%   \begin{equation}\label{eq:5}
%     Z_{22} = \frac{V_2}{I_2}\Bigr|_{\substack{I_1=0}}
%   \end{equation}
% \end{multicols}

% \begin{itemize}
%     \item {Equação \ref{eq:2}: Impedância de entrada de circuito aberto;}
%     \item {Equação \ref{eq:3}: Impedância de transferência de circuito aberto da porta 1 para a porta 2;}
%     \item {Equação \ref{eq:4}: Impedância de transferência de circuito aberto da porta 2 para a porta 1}
%     \item {Equação \ref{eq:5}: Impedância de saída de circuito aberto;}
% \end{itemize}

% Uma vez que são determinados os parâmetros de impedância, é fácil obter os parâmetros de admitância através da aplicação das seguintes fórmulas de transformação\cite{b1}.

% \begin{multicols}{2}
%   \begin{equation}\label{eq:6}
%     Y_{11} = \frac{Z_{22}}{\Delta_z}
%   \end{equation}\break
%   \begin{equation}\label{eq:7}
%     Y_{12} = -\frac{Z_{21}}{\Delta_z}
%   \end{equation}
% \end{multicols}

% \begin{multicols}{2}
%   \begin{equation}\label{eq:8}
%     Y_{21} = -\frac{Z_{12}}{\Delta_z}
%   \end{equation}\break
%   \begin{equation}\label{eq:9}
%     Y_{22} = \frac{Z_{11}}{\Delta_z}
%   \end{equation}
% \end{multicols}